\documentclass[aps,prd,onecolumn,longbibliography,nofootinbib]{revtex4-2}

\usepackage[utf8]{inputenc}
\usepackage{amsmath,amssymb,mathtools}
\usepackage{graphicx}
\usepackage{xcolor}
\usepackage{hyperref}
\usepackage{enumitem}
\usepackage{booktabs}
\usepackage[final]{microtype}

\hypersetup{colorlinks=true, linkcolor=blue, citecolor=blue, urlcolor=blue}

\newcommand{\Tr}{\mathrm{Tr}}
\newcommand{\Hcore}{\mathcal{H}_{\rm core}}
\newcommand{\Hrad}{\mathcal{H}_{\rm rad}}
\newcommand{\Htot}{\mathcal{H}_{\rm tot}}
\newcommand{\Smicro}{S_{\rm micro}}
\newcommand{\Svn}{S_{\rm vN}}

\begin{document}

\title{Thermal Radiation and Accelerating Evaporation in a Non-Gravitational Sector Model}

\author{[Author]}
\affiliation{[Affiliation]}
\date{\today}

\begin{abstract}
We introduce a finite sector model for studying thermodynamic and
information-theoretic features of black-hole evaporation without dynamical
gravity.  The model is governed by transition rates derived from a Hamiltonian
coupling between adjacent core sectors and outgoing radiation modes.  The core
sectors have an area-like state count, and a square-root relation between
sector label and energy gives the Schwarzschild scalings $S\sim M^2$,
$T\sim M^{-1}$, and negative heat capacity.  In an energy-resolved
implementation, the microcanonical density-of-states ratio gives a hard
spectrum close to $x^p e^{-x}$, with $x=\beta\omega$.  With intra-sector
equilibration, the same dynamics gives accelerating emission for the
square-root mass law, while a linear mass-law control remains approximately
non-accelerating despite producing a similar hard spectrum.  We also study a
finite unitary record model with the same shrinking-support architecture and
find a Page-like radiation-entropy turnover and enhanced early/late radiation
mutual information in scrambled runs.  The results identify a finite
quantum-statistical mechanism for thermal spectra, acceleration, and unitary
information flow in black-hole-like evaporation models.
\end{abstract}

\maketitle

\section{Introduction}

Black-hole evaporation combines thermodynamic and quantum-information
statements.  A four-dimensional Schwarzschild black hole has entropy
proportional to area, temperature inversely proportional to mass, negative heat
capacity, and power scaling $P\sim M^{-2}$ \cite{Bekenstein1973,Hawking1975}.
Page's argument gives the expected fine-grained radiation entropy for unitary
evaporation \cite{Page1993Info,Page1993Average}.  Modern gravitational
calculations recover this behavior using generalized entropy, quantum extremal
surfaces, islands, and replica-wormhole saddles
\cite{Penington2020,AEMM2019,AlmheiriReview2021}.

This paper asks how much of that package can be reproduced in a finite quantum
model without dynamical spacetime.  The point is comparative.  A model that
uses only ordinary Hilbert spaces, sector transitions, density-of-states
effects, and unitary purification can help separate generic statistical
mechanisms from specifically gravitational ones.

Toy and qubit models of black-hole evaporation already show that unitary
information transfer and Page-like behavior can be modeled without solving the
gravitational problem \cite{Avery2013,OsugaPage2018,Broda2021}.  Tensor-network
and random-unitary models define explicit maps between black-hole and radiation
Hilbert spaces \cite{LeutheusserVanRaamsdonk2017,PiroliSunderhaufQi2020}.  The
present construction follows that explicit-model tradition, with a different
emphasis: the thermodynamic exponents, the hard spectrum, and the information
diagnostics are kept as separate tests.

The central model is a sector-rate model derived from an abstract Hamiltonian.
Core sectors $B_n$ have dimensions
\begin{equation}
    \dim B_n=q^n,\qquad \Smicro(n)=n\log q,
\end{equation}
and the black-hole-like mass law is
\begin{equation}
    M_n=\alpha\sqrt n .
\end{equation}
This gives $S\sim M^2$, $T\sim M^{-1}$, and negative heat capacity.  The
comparison mass law
\begin{equation}
    M_n=\alpha n
\end{equation}
keeps the same sector count and removes the black-hole-like entropy/mass
relation.

The main numerical result is that an energy-resolved sector model can generate
both a near-thermal hard spectrum and accelerated evaporation, provided the
remaining core is mixed within each sector.  In the same sector model, the
linear mass-law control gives a similarly good hard spectrum and does not show
the same acceleration.  This isolates the acceleration as a consequence of the
entropy/mass relation rather than of the radiation alphabet alone.

The information-flow calculation is separate.  It is a small exact
state-vector repeated-interaction realization designed to track purification,
soft records, and early/late radiation correlations.  It should be read as a
compatibility check: the sector thermodynamics can coexist with a unitary
record model showing a Page-like turnover.  It is not yet a full coherent
Hamiltonian simulation of the sector-rate model.

\section{Model Standard and Entropies}
\label{sec:standard}

The minimum object is:
\begin{enumerate}[leftmargin=*]
    \item a finite Hilbert space;
    \item a Hamiltonian or Hamiltonian-derived transition rule;
    \item an initial state class;
    \item an observable split into core, hard radiation, purifier or bath, and
    full emitted radiation;
    \item controls that change the mass law, density of states, and mixing.
\end{enumerate}

The microcanonical core entropy is
\begin{equation}
    \Smicro(n)=\log\Omega_n,\qquad \Omega_n=\dim B_n .
\end{equation}
It is the entropy entering the thermodynamic relation
\begin{equation}
    \beta=\frac{\partial \Smicro}{\partial M}.
\end{equation}
The hard-radiation entropy is the entropy of the visible hard register or hard
energy distribution.  The Page-like diagnostic is the von Neumann entropy of
the full emitted radiation subsystem,
\begin{equation}
    \Svn(R)=-\Tr \rho_R\log\rho_R .
\end{equation}
These quantities answer different questions and are reported separately.

\section{Sector Hamiltonian and Rate Limit}
\label{sec:sector}

The abstract total Hilbert space is
\begin{equation}
    \Htot=
    \left(\bigoplus_{n=n_{\min}}^{n_{\max}} B_n\right)\otimes \Hrad .
\end{equation}
The Hamiltonian has the form
\begin{equation}
    H_{\rm tot}=H_{\rm core}+H_{\rm rad}+H_{\rm int}.
\end{equation}
The core block is
\begin{equation}
    H_{\rm core}=\sum_{n,a}E_{n,a}|n,a\rangle\langle n,a|,
    \qquad
    E_{n,a}=M_n+\epsilon_{n,a}.
\end{equation}
The radiation Hamiltonian is
\begin{equation}
    H_{\rm rad}=\sum_\lambda \omega_\lambda b_\lambda^\dagger b_\lambda .
\end{equation}
The interaction couples one core sector to the next smaller sector plus one
outgoing quantum:
\begin{equation}
    H_{\rm int}
    =
    \sum_{n,a,b,\lambda}
    \left(
    g^{(\lambda)}_{n;ba}|n-1,b\rangle\langle n,a|b_\lambda^\dagger
    +{\rm h.c.}
    \right).
\end{equation}
In the weak-coupling secular limit,
\begin{equation}
    \Gamma_{n,a\rightarrow n-1,b;\lambda}
    =
    2\pi |g^{(\lambda)}_{n;ba}|^2
    \delta(E_{n,a}-E_{n-1,b}-\omega_\lambda).
\end{equation}
After binning or continuum approximation, this is a golden-rule rate: matrix
element times radiation phase space times final core density of states.

The population diagnostic evolves probabilities by
\begin{equation}
    \dot p_{n,a}
    =
    -\sum_b \Gamma_{ba}p_{n,a}
    +\sum_c \Gamma_{ac}p_{n+1,c}.
\end{equation}
This rate equation is the object actually simulated in the sector
calculation.  A fully coherent wave-packet simulation remains a later
Hamiltonian upgrade.

\section{Thermodynamics}
\label{sec:thermo}

For the square-root mass law,
\begin{equation}
    M_n=\alpha\sqrt n,\qquad \Smicro(n)=n\log q .
\end{equation}
Therefore
\begin{equation}
    \beta_n=
    \frac{\partial \Smicro}{\partial M}
    =
    \frac{2\sqrt n\log q}{\alpha},
    \qquad
    T_n=\frac{\alpha}{2\sqrt n\log q}.
\end{equation}
The heat capacity is negative because $M$ decreases as $T$ increases.  This is
a consequence of the chosen entropy/mass relation, not an independent
dynamical discovery.

For the linear mass law,
\begin{equation}
    M_n=\alpha n ,
\end{equation}
the temperature is constant for the same sector count.  This gives the main
thermodynamic control.

For a transition emitting energy $\omega$, the microcanonical final-state
factor is
\begin{equation}
    \exp[\Smicro(M-\omega)-\Smicro(M)].
\end{equation}
For small $\omega$,
\begin{equation}
    \Smicro(M-\omega)-\Smicro(M)
    =
    -\beta(M)\omega+O(\omega^2).
\end{equation}
Thus a thermal hard spectrum can arise from sector density of states and
transition matrix elements.  This is the thermality test used below.  It is
different from assigning a hard-register probability and then purifying it.

\section{Energy-Resolved Sector Result}
\label{sec:energy}

The first sector-rate diagnostic used small random perturbations inside each
sector.  It produced accelerating evaporation for the square-root mass law and
deceleration for the linear mass law.  Its emitted spectrum was too narrow,
with total-variation distance from the target thermal distribution around
$0.77$.

The improved diagnostic gives each sector an internal spectrum
\begin{equation}
    E_{n,a}=M_n+\epsilon_{n,a}
\end{equation}
with local density of states approximately
\begin{equation}
    \rho_n(\epsilon)\propto e^{\beta_n\epsilon}
\end{equation}
over an energy window of order $T_n$.  Transition rates are
\begin{equation}
    \Gamma_{fi}\propto
    |\langle f,n-1|X_n|i,n\rangle|^2\,\omega_{fi}^{p},
    \qquad
    \omega_{fi}=E_i^{(n)}-E_f^{(n-1)}.
\end{equation}
The emitted number distribution is compared in
\begin{equation}
    x=\beta_n\omega
\end{equation}
to
\begin{equation}
    P(x)\propto x^p e^{-x}.
\end{equation}

With an exponential sector density of states and no within-sector mixing, the
hard spectrum becomes close to thermal and the evaporation decelerates.  The
state falls into lower-rate parts of the sector spectrum.  The next diagnostic
adds an idealized within-sector mixing step,
\begin{equation}
    p_n(a)\mapsto \frac{P_n}{\dim B_n}.
\end{equation}
This stands in for fast intra-sector equilibration.

With this mixing, the square-root mass law gives both acceleration and a
near-thermal hard spectrum.  For two seeds, local and scrambled shrinkage
operators, and width parameter ${\rm width}_x=4$:
\begin{center}
\begin{tabular}{lccc}
\toprule
operator & power ratio & jump ratio & TV distance\\
\midrule
local     & 1.098 & 1.045 & 0.055\\
scrambled & 1.097 & 1.045 & 0.054\\
\bottomrule
\end{tabular}
\end{center}
The power ratio is mid-evaporation emitted power divided by early emitted
power.

The linear mass-law control has a similarly good hard spectrum and lacks the
same acceleration:
\begin{center}
\begin{tabular}{lccc}
\toprule
operator & power ratio & jump ratio & TV distance\\
\midrule
local     & 0.988 & 0.995 & 0.051\\
scrambled & 0.988 & 0.995 & 0.049\\
\bottomrule
\end{tabular}
\end{center}

This is the main sector-rate result.  The hard spectrum is generated by the
energy-resolved density of states.  Acceleration is tied to the square-root
mass law.  Within-sector mixing is still a modeled ingredient.

\section{Companion Information-Flow Diagnostic}
\label{sec:info}

The sector-rate calculation does not track global purification.  The companion
diagnostic is a finite repeated-interaction state-vector model with core
sectors
\begin{equation}
    \dim B_L=q^{L^2}.
\end{equation}
It includes hard radiation records, hidden purifier records, soft shell
records, and an emitted-energy accumulator.  The hard marginal in this
companion model is fixed by the chosen thermodynamic schedule.  The
nontrivial question is whether the shrinking support and soft records can
carry unitary information flow in the same finite architecture.

The main small run has
\begin{equation}
    L_{\rm init}=3,\qquad q=2,
\end{equation}
and uses one fixed eight-emission hard schedule with threshold $\Delta=4$.
The following table reads off the same trajectory at four times:
\begin{center}
\begin{tabular}{lcccc}
\toprule
quantity & 2 & 4 & 6 & 8\\
\midrule
scrambled $\Svn(R_{\rm full})$ & 1.427 & 3.799 & 2.612 & 1.900\\
no-scrambling $\Svn(R_{\rm full})$ & 1.017 & 1.384 & 1.592 & 1.698\\
$\log|\mathrm{supp}({\rm core},A)|$ & 6.947 & 4.159 & 3.584 & 2.197\\
mean active capacity & 5.666 & 2.716 & 1.154 & 0.401\\
\bottomrule
\end{tabular}
\end{center}
The support and capacity rows are the same for the two scramblers on this
schedule.  The scrambled run shows a small rise and turnover of the full
radiation von Neumann entropy.  The no-scrambling run rises more slowly over
the same window.  This indicates that the finite turnover observed here is not
only a dimension-bound effect; routing information into the emitted records
also matters.

For the six-emission old/new split, the full radiation mutual information is:
\begin{center}
\begin{tabular}{lc}
\toprule
scrambler & $I({\rm old}:{\rm new})$\\
\midrule
Margulis-like & 2.615\\
grid-like     & 2.624\\
none          & 1.988\\
\bottomrule
\end{tabular}
\end{center}
This is a small-system compatibility check.  It does not prove a robust
thermodynamic-limit Page curve for the sector-rate model.

\section{Inputs, Outputs, and Controls}
\label{sec:accounting}

The specified inputs are:
\begin{itemize}[leftmargin=*]
    \item the sector dimensions $\dim B_n=q^n$;
    \item the square-root and linear mass laws;
    \item the energy-resolved sector density of states;
    \item the transition operators $X_n$;
    \item the idealized within-sector mixing map;
    \item the purifier and soft-record registers in the companion model.
\end{itemize}

The generated outputs in the sector-rate diagnostic are:
\begin{itemize}[leftmargin=*]
    \item hard spectrum close to $x^p e^{-x}$ for exponential local density of
    states;
    \item acceleration for the square-root mass law after within-sector mixing;
    \item absence of the same acceleration for the linear mass-law control;
    \item failure of broad-spectrum thermality alone to preserve acceleration
    without mixing.
\end{itemize}

The generated outputs in the companion information-flow diagnostic are:
\begin{itemize}[leftmargin=*]
    \item a fixed-schedule radiation entropy rise and turnover in the
    scrambled run;
    \item a shrinking core-support bound on the same trajectory;
    \item enhanced early/late mutual information with scrambling.
\end{itemize}

The main remaining gap is that these are two linked diagnostics, not one
coherent simulation.  A stronger model would put the sector-rate dynamics,
explicit outgoing modes, mixing dynamics, and information-flow diagnostics in
one state-vector or open-system Hamiltonian calculation.

\section{Toward a Microscopic Origin}
\label{sec:microscopic}

The sector-rate model specifies the area-like state count and mass law.  A
natural microscopic derivation would be stronger.  The current live direction
is a relational connector model: a core with $N$ active constituents and
$O(N^2)$ pairwise connector modes.  Such a model can supply area-like counting
from ordinary degrees of freedom.

The spectral constraints are now sharp.  A linear critical connector band gives
a good simple power proxy and poor post-emission heating.  A quadratic
($z=2$) connector band gives better heating.  For that quadratic band, the
emission matrix element must scale approximately as
\begin{equation}
    |g(\omega)|^2\sim \omega^{1/2}
\end{equation}
to recover $P\sim M^{-2}$ in the local spectral-weight proxy.  This points
toward type-B Goldstone or ferromagnetic-magnon-like connector physics.

The open microscopic problem is to obtain, from one recognizable Hamiltonian,
the $O(N^2)$ active connector count, protected quadratic band, emission
matrix-element scaling, and slow shrinkage after many microscopic emissions.

\section{Discussion}

The sector-rate construction gives a reproducible non-gravitational comparison
model.  It assembles an area-like sector count, black-hole-like mass law,
Hamiltonian-derived transition rates, a generated near-thermal hard spectrum,
and acceleration under the square-root mass law.  The companion
repeated-interaction diagnostic shows compatible unitary information flow at
small finite size.

The claim is limited.  The area state count, mass law, sector density of states,
and within-sector mixing are specified.  The information-flow diagnostic is a
companion construction rather than the same sector-rate Hamiltonian evolved
coherently.  The finite Page-like signal is small and seed-limited.  The
fully coherent autonomous Hamiltonian with outgoing radiation wave packets
remains to be built.

Within those limits, the model helps separate the evaporation package into
testable mechanisms.  Thermality comes from energy-resolved density of states;
acceleration comes from the entropy/mass relation plus mixing; Page-like
information flow requires a shrinking support and actual routing of
correlations into emitted records.  Those statements are ordinary finite
quantum statistical mechanics.  A gravitational derivation still has to supply
the sector count, the mass law, and the microscopic dynamics that realize them.

\appendix

\section{Reproducibility Notes}

The main scripts and notes are:
\begin{center}
\begin{tabular}{ll}
\toprule
script & result note or output\\
\midrule
\texttt{sim/sector\_hamiltonian\_evaporator.py} &
\texttt{notes/sector\_hamiltonian\_evaporator\_results.md}\\
\texttt{sim/sector\_hamiltonian\_energy\_resolved.py} &
\texttt{notes/sector\_hamiltonian\_energy\_resolved\_results.md}\\
\texttt{sim/repeated\_interaction\_page\_trajectory.py} &
\texttt{sim/data/repeated\_interaction\_page\_trajectory.csv}\\
\texttt{sim/fused\_floquet\_time\_resolved\_scan.py} &
\texttt{notes/fused\_floquet\_time\_resolved\_results.md}\\
\bottomrule
\end{tabular}
\end{center}

The energy-resolved sector runs use
\begin{equation}
    n=4,\ldots,10,\qquad q=2,\qquad \alpha=8,
\end{equation}
two seeds, local and scrambled shrinkage maps, and exponential sector density
of states with width parameters ${\rm width}_x=2,4,8$.  The displayed tables
use ${\rm width}_x=4$.

The companion information-flow run uses $L_{\rm init}=3$, $q=2$, threshold
$\Delta=4$, one fixed eight-emission hard schedule, and the Margulis-like and
no-scrambling choices.  The full radiation subsystem includes hard records,
hidden purifier records, and soft shell records.

\bibliographystyle{apsrev4-2}
\bibliography{refs}

\end{document}
